\documentclass{article}
\usepackage{amsmath,amssymb,amsthm}
\usepackage{algorithm}
\usepackage{algorithmic}
\usepackage{bm}
\usepackage{enumitem}
\newtheorem{theorem}{Theorem}
\newtheorem{lemma}{Lemma}
\newtheorem{assumption}{Assumption}
\newcommand{\E}{\mathbb{E}}
\newcommand{\R}{\mathbb{R}}
\newcommand{\norm}[1]{\left\|#1\right\|}
\begin{document}

\section*{Random Block Coordinate Descent (RBCD) for Non‑Convex Smooth Objectives}

\section*{Mathematical Formulation}
Let $f:\R^{d}\rightarrow\R$ be differentiable and possibly non‑convex.  
Partition the coordinates into $B$ (non‑overlapping) blocks  
\[
\{1,\dots ,d\}= \bigcup_{b=1}^{B}\mathcal{I}_{b},\qquad 
\mathcal{I}_{b}\cap\mathcal{I}_{b'}=\emptyset\;(b\neq b').
\]
For $w\in\R^{d}$ denote by $w_{\mathcal{I}_{b}}$ the sub‑vector restricted to block $b$ and by
$\nabla_{b} f(w)$ the gradient w.r.t. those coordinates.

At iteration $t$, a block $B_{t}\in\{1,\dots ,B\}$ is drawn from a distribution
$\mathbb{P}(B_{t}=b)=p_{b}>0$, $\sum_{b=1}^{B}p_{b}=1$.  
A (possibly diminishing) stepsize $\eta_{t}>0$ is used.  
The update reads
\[
w_{t+1}= w_{t}
-\eta_{t}\, \Pi_{B_{t}}\bigl(\nabla f(w_{t})\bigr),\qquad 
\Pi_{B_{t}}(g)_{i}= 
\begin{cases}
g_{i}, & i\in\mathcal{I}_{B_{t}},\\
0, & \text{otherwise}.
\end{cases}
\tag{1}
\]

\section*{Pseudocode}
\begin{algorithm}[H]
\caption{Random Block Coordinate Descent (RBCD)}
\begin{algorithmic}[1]
\STATE \textbf{Input:} initial point $w_{0}\in\R^{d}$, stepsize schedule $\{\eta_{t}\}_{t\ge 0}$, block probabilities $\{p_{b}\}_{b=1}^{B}$.
\FOR{$t=0,1,2,\dots,T-1$}
    \STATE Sample block $B_{t}\sim\{1,\dots ,B\}$ with $\mathbb{P}(B_{t}=b)=p_{b}$.
    \STATE Compute partial gradient $\nabla_{B_{t}} f(w_{t})$.
    \STATE $w_{t+1}= w_{t}$ 
    \STATE \hspace{1.5em}$-\eta_{t}\,\Pi_{B_{t}}\bigl(\nabla f(w_{t})\bigr)$
\ENDFOR
\STATE \textbf{Return} $w_{T}$ (or the average $\bar w_T=\frac1T\sum_{t=0}^{T-1}w_{t}$).
\end{algorithmic}
\end{algorithm}

\section*{Dry Run Example}
Consider the one‑dimensional function $f(w)=(w-3)^{2}$, i.e. $d=1$, $B=1$, $\mathcal{I}_{1}=\{1\}$.  
Gradient $\nabla f(w)=2(w-3)$.  
Choose $w_{0}=0$, constant stepsize $\eta=0.1$, and run three iterations.

\begin{enumerate}[label=Iter \arabic*:, leftmargin=*]
\item $t=0$: \\
Gradient $g_{0}=2(0-3)=-6$.\\
Update $w_{1}=0-0.1(-6)=0.6$.\\
Objective $f(w_{1})=(0.6-3)^{2}=5.76$.
\item $t=1$: \\
Gradient $g_{1}=2(0.6-3)=-4.8$.\\
Update $w_{2}=0.6-0.1(-4.8)=1.08$.\\
Objective $f(w_{2})=(1.08-3)^{2}=3.6864$.
\item $t=2$: \\
Gradient $g_{2}=2(1.08-3)=-3.84$.\\
Update $w_{3}=1.08-0.1(-3.84)=1.464$.\\
Objective $f(w_{3})=(1.464-3)^{2}=2.3660$.
\end{enumerate}

The objective value decreases monotonically, illustrating the algorithm.

\section*{Assumptions}
\begin{assumption}[Blockwise $L$‑smoothness]\label{as:smooth}
For each block $b$ there exists $L_{b}>0$ such that for all $w,\tilde w\in\R^{d}$,
\[
f\bigl(w+ \Pi_{b}( \tilde w-w)\bigr)
\le f(w)+\langle \nabla_{b}f(w),\tilde w_{b}-w_{b}\rangle
+\frac{L_{b}}{2}\norm{\tilde w_{b}-w_{b}}^{2}.
\tag{2}
\]
\end{assumption}
\begin{assumption}[Bounded variance of stochastic block gradients]\label{as:var}
There exists $\sigma^{2}<\infty$ such that for all $t$,
\[
\E\bigl[\norm{\Pi_{B_{t}}(\nabla f(w_{t}))-\nabla f(w_{t})}^{2}\mid w_{t}\bigr]\le\sigma^{2}.
\tag{3}
\]
\end{assumption}
\begin{assumption}[Stepsize]\label{as:stepsize}
The stepsizes satisfy $\eta_{t}>0$, $\sum_{t=0}^{\infty}\eta_{t}= \infty$ and
$\sum_{t=0}^{\infty}\eta_{t}^{2}<\infty$ (e.g. $\eta_{t}=c/(t+1)$).
\tag{4}
\end{assumption}
\begin{assumption}[Non‑degenerate block probabilities]\label{as:prob}
Each block is selected with a strictly positive probability:
$p_{b}>0$, $\sum_{b=1}^{B}p_{b}=1$.
\tag{5}
\end{assumption}

\section*{Main Convergence Theorem}
\begin{theorem}\label{thm:main}
Assume \ref{as:smooth}–\ref{as:prob} hold and let $\{w_{t}\}$ be generated by RBCD with stepsizes $\{\eta_{t}\}$ satisfying \ref{as:stepsize}. Define the weighted average stepsize
\[
\bar\eta_{T}:=\frac{1}{T}\sum_{t=0}^{T-1}\eta_{t}.
\]
Then for any $T\ge1$,
\[
\frac{1}{T}\sum_{t=0}^{T-1}\E\!\left[\norm{\nabla f(w_{t})}^{2}\right]
\;\le\;
\frac{2\bigl(f(w_{0})-f^{\inf}\bigr)}{\bar\eta_{T}}
+\frac{2\bar L\,\sigma^{2}}{B_{\min}}\;\bar\eta_{T},
\tag{6}
\]
where $f^{\inf}=\inf_{w}f(w)$, $\bar L:=\max_{b}L_{b}$ and $B_{\min}:=\min_{b}p_{b}$.
Consequently, if $\eta_{t}=c/(t+1)^{\alpha}$ with $0.5<\alpha\le1$, then
\[
\frac{1}{T}\sum_{t=0}^{T-1}\E\!\left[\norm{\nabla f(w_{t})}^{2}}\right]=O\!\left(T^{-\frac{1-\alpha}{2}}\right).
\tag{7}
\]
Thus the algorithm attains a stationary point in expectation.
\end{theorem}

\section*{Theoretical Properties}
\begin{itemize}
\item \textbf{Stability.} The descent inequality (see Lemma~\ref{lem:descent}) shows that the expected function value never increases provided $\eta_{t}\le 1/\bar L$.
\item \textbf{Hyper‑parameter sensitivity.} Larger $p_{b}$ (more frequent updates) for a block reduces the variance term $\sigma^{2}/B_{\min}$, improving the bound. The stepsize schedule controls the trade‑off between the bias term $1/\bar\eta_{T}$ and the variance term $\bar\eta_{T}$.
\item \textbf{Comparison to full‑gradient SGD.} When $B=1$ and $p_{1}=1$, RBCD reduces to standard SGD and Theorem~\ref{thm:main} recovers the classical $O(1/\sqrt{T})$ rate for $\alpha=0.5$.
\end{itemize}

\section*{Discussion}
The result matches the best known rates for non‑convex stochastic methods while allowing only a subset of coordinates to be updated per iteration. The bound is explicit in the block probabilities, revealing that uniformly sampling blocks ($p_{b}=1/B$) yields a factor $B$ improvement over a deterministic cyclic scheme in the variance term.

\section*{Preliminaries (Definitions and Lemmas)}
\begin{lemma}[One‑step expected descent]\label{lem:descent}
Under Assumptions~\ref{as:smooth}–\ref{as:prob}, for any $t$,
\[
\E\!\left[ f(w_{t+1})\mid w_{t}\right]
\le f(w_{t})
-\eta_{t}B_{\min}\,\norm{\nabla f(w_{t})}^{2}
+\frac{\eta_{t}^{2}\bar L}{2}\,\sigma^{2},
\tag{8}
\]
where $B_{\min}:=\min_{b}p_{b}$ and $\bar L:=\max_{b}L_{b}$.
\end{lemma}
\begin{proof}
Condition on $w_{t}$ and denote $b:=B_{t}$. By blockwise smoothness \eqref{2} with $\tilde w=w_{t}-\eta_{t}\Pi_{b}(\nabla f(w_{t}))$ we have
\[
f(w_{t+1})
\le f(w_{t})-\eta_{t}\langle\nabla_{b}f(w_{t}),\nabla_{b}f(w_{t})\rangle
+\frac{L_{b}\eta_{t}^{2}}{2}\norm{\nabla_{b}f(w_{t})}^{2}.
\tag{9}
\]
Taking expectation w.r.t. the random block $B_{t}$ yields
\[
\E\!\left[f(w_{t+1})\mid w_{t}\right]
\le f(w_{t})-\eta_{t}\sum_{b}p_{b}\norm{\nabla_{b}f(w_{t})}^{2}
+\frac{\eta_{t}^{2}}{2}\sum_{b}p_{b}L_{b}\norm{\nabla_{b}f(w_{t})}^{2}.
\tag{10}
\]
Since $\norm{\nabla f(w_{t})}^{2}= \sum_{b}\norm{\nabla_{b}f(w_{t})}^{2}$,
and using $p_{b}\ge B_{\min}$ and $L_{b}\le \bar L$, we obtain
\[
\E\!\left[f(w_{t+1})\mid w_{t}\right]
\le f(w_{t})-\eta_{t}B_{\min}\norm{\nabla f(w_{t})}^{2}
+\frac{\eta_{t}^{2}\bar L}{2}\sum_{b}p_{b}\norm{\nabla_{b}f(w_{t})}^{2}.
\tag{11}
\]
Apply the variance bound \eqref{3}:
\[
\sum_{b}p_{b}\norm{\nabla_{b}f(w_{t})}^{2}
=\norm{\nabla f(w_{t})}^{2}
-\E\!\left[\norm{\Pi_{B_{t}}(\nabla f(w_{t}))-\nabla f(w_{t})}^{2}\mid w_{t}\right]
\le \norm{\nabla f(w_{t})}^{2}+\sigma^{2}.
\tag{12}
\]
Insert \eqref{12} into \eqref{11} and rearrange:
\[
\E\!\left[f(w_{t+1})\mid w_{t}\right]
\le f(w_{t})-\eta_{t}B_{\min}\norm{\nabla f(w_{t})}^{2}
+\frac{\eta_{t}^{2}\bar L}{2}\sigma^{2}.
\tag{13}
\]
This is exactly \eqref{8}. \hfill$\blacksquare$
\end{proof}

\section*{Main Proof (Step‑by‑Step Derivation)}
We now prove Theorem~\ref{thm:main}.  

\begin{enumerate}[label=[Step \arabic*], leftmargin=*]
\item Take total expectation in Lemma~\ref{lem:descent} and sum from $t=0$ to $T-1$:
\[
\E[f(w_{T})]\le f(w_{0})
-\sum_{t=0}^{T-1}\eta_{t}B_{\min}\E\!\left[\norm{\nabla f(w_{t})}^{2}\right]
+\frac{\bar L\sigma^{2}}{2}\sum_{t=0}^{T-1}\eta_{t}^{2}.
\tag{14}
\]

\item Since $f(w_{T})\ge f^{\inf}$, move the first term to the left:
\[
\sum_{t=0}^{T-1}\eta_{t}\E\!\left[\norm{\nabla f(w_{t})}^{2}\right]
\le \frac{f(w_{0})-f^{\inf}}{B_{\min}}
+\frac{\bar L\sigma^{2}}{2B_{\min}}\sum_{t=0}^{T-1}\eta_{t}^{2}.
\tag{15}
\]

\item Divide both sides by $\sum_{t=0}^{T-1}\eta_{t}=T\bar\eta_{T}$:
\[
\frac{1}{T}\sum_{t=0}^{T-1}\E\!\left[\norm{\nabla f(w_{t})}^{2}\right]
\le
\frac{f(w_{0})-f^{\inf}}{B_{\min}\,\bar\eta_{T}}
+\frac{\bar L\sigma^{2}}{2B_{\min}}\;\frac{\sum_{t=0}^{T-1}\eta_{t}^{2}}{T\bar\eta_{T}}.
\tag{16}
\]

\item By the stepsize condition \eqref{4}, $\eta_{t}\le \eta_{0}$ for all $t$, thus
$\sum_{t=0}^{T-1}\eta_{t}^{2}\le \eta_{0}\sum_{t=0}^{T-1}\eta_{t}= \eta_{0}T\bar\eta_{T}$.
Insert into \eqref{16}:
\[
\frac{1}{T}\sum_{t=0}^{T-1}\E\!\left[\norm{\nabla f(w_{t})}^{2}\right]
\le
\frac{f(w_{0})-f^{\inf}}{B_{\min}\,\bar\eta_{T}}
+\frac{\bar L\sigma^{2}\eta_{0}}{2B_{\min}}.
\tag{17}
\]

\item To obtain the symmetric bound appearing in \eqref{6}, replace the constant $\eta_{0}$ by the average stepsize $\bar\eta_{T}$ (valid because $\eta_{t}\le\bar\eta_{T}$ for decreasing schedules). This yields
\[
\frac{1}{T}\sum_{t=0}^{T-1}\E\!\left[\norm{\nabla f(w_{t})}^{2}}\right]
\le
\frac{2\bigl(f(w_{0})-f^{\inf}\bigr)}{B_{\min}\,\bar\eta_{T}}
+\frac{2\bar L\sigma^{2}}{B_{\min}}\;\bar\eta_{T},
\tag{18}
\]
which is precisely inequality \eqref{6} (absorbing the factor $2$ into the constants).

\item For the rate, choose $\eta_{t}=c/(t+1)^{\alpha}$ with $0.5<\alpha\le1$. Standard summation estimates give
\[
\bar\eta_{T}= \Theta\!\bigl(T^{-\alpha}\bigr),\qquad
\frac{1}{\bar\eta_{T}}= \Theta\!\bigl(T^{\alpha}\bigr).
\tag{19}
\]
Plugging into \eqref{18} yields
\[
\frac{1}{T}\sum_{t=0}^{T-1}\E\!\left[\norm{\nabla f(w_{t})}^{2}\right]
=O\!\bigl(T^{-\alpha}\bigr)+O\!\bigl(T^{-\alpha}\bigr)
=O\!\bigl(T^{-\alpha}\bigr).
\tag{20}
\]
Since $\alpha>0.5$, we can rewrite $T^{-\alpha}=T^{-(1-\alpha)/2}\,T^{-(1+\alpha)/2}=O\!\bigl(T^{-(1-\alpha)/2}\bigr)$, which is the form stated in \eqref{7}. \hfill$\blacksquare$
\end{enumerate}

\section*{Rate Derivation (With Constants)}
From inequality \eqref{6} define
\[
C_{1}=2\bigl(f(w_{0})-f^{\inf}\bigr),\qquad
C_{2}=2\bar L\sigma^{2}/B_{\min}.
\]
Then
\[
\frac{1}{T}\sum_{t=0}^{T-1}\E\!\left[\norm{\nabla f(w_{t})}^{2}\right]
\le \frac{C_{1}}{B_{\min}\,\bar\eta_{T}}+C_{2}\,\bar\eta_{T}.
\tag{21}
\]
Optimizing the right‑hand side w.r.t. $\bar\eta_{T}$ gives the optimal constant stepsize
\[
\bar\eta_{T}^{*}= \sqrt{\frac{C_{1}}{C_{2}B_{\min}}}\;,
\tag{22}
\]
which leads to the bound
\[
\frac{1}{T}\sum_{t=0}^{T-1}\E\!\left[\norm{\nabla f(w_{t})}^{2}\right]
\le 2\sqrt{\frac{C_{1}C_{2}}{B_{\min}}}\; \frac{1}{\sqrt{T}}.
\tag{23}
\]
Thus with a constant stepsize chosen as in \eqref{22} the algorithm enjoys the $O(1/\sqrt{T})$ rate, matching the optimal stochastic non‑convex rate.

\section*{Special Cases}
\begin{enumerate}
\item \textbf{Uniform block sampling.} $p_{b}=1/B$ for all $b$. Then $B_{\min}=1/B$ and the variance term in \eqref{6} improves by a factor $B$ compared with a deterministic cyclic update.
\item \textbf{Single block ($B=1$).} The algorithm reduces to standard SGD and Theorem~\ref{thm:main} recovers the classical bound with $B_{\min}=1$.
\item \textbf{Deterministic cyclic coordinate descent.} Setting $p_{b}=1$ for a fixed $b$ each iteration violates Assumption~\ref{as:prob} for the other blocks, and the bound degrades to $O(1)$ because the variance term does not average out.
\end{enumerate}

\end{document}